\chapter{LITERATURE REVIEW}


\newcommand{\funcgroup}[1]{%
	\vspace{\baselineskip}%
	\noindent{\bfseries\fontsize{12}{14.4}\selectfont #1}%
	\vspace{0.1pt}\par%
}

\begin{comment}
	%for counter
	\newcounter{funcgroupcounter}[section]  % resets to 0 each new \section
	\newcommand{\funcgroup}[1]{%
		\refstepcounter{funcgroupcounter}%
		\vspace{\baselineskip}%
		\noindent{\bfseries\fontsize{12}{14.4}\selectfont
			\thesection.\arabic{funcgroupcounter}\quad #1}%
		\vspace{0.1pt}\par%
	}
\end{comment}

This chapter discusses the published academic papers used to support our project. There are two main parts: The first part discussed the technology stack, website development framework, and website development planning that has been implemented and met with success. Such technology stack and planning would then be used as a guideline to our project development. The second part explores any methods or algorithms that could be adapted and used for the project’s main functions development. For this part, it will be divided into 3 sub-parts for 3 functions, with each part containing reviewed literature that are related to the function of that part.

\section{Technological Stack and System Architecture}

\subsection{Design web service academic information system based multiplatform \cite{paper_10}}

This paper proposed a system design that uses Web Services to make academic information available on both web and mobile platforms.

The core of the design is a REST API that uses JSON to exchange data between the server and different types of devices. This approach was chosen because it allows a single database to support many different user interfaces while keeping all the information synchronized in real-time.

The study found that using this web service architecture makes the system much easier to maintain and scale up because the "business logic" on the server is kept separate from how the app looks on the screen.

For future work, the authors suggested adding stronger security protocols to the API and making the database faster so it can handle more students using the app at the same time.


\subsection{Design and Implementation of a Digital Notice Board and Forum Discussion System \cite{paper_11}}

This paper presents an educational website, equipped with two functions, digital notice board and discussion forum, which are developed to fulfill two goals that aim to improve university students' quality of life. The first goal is to create a centralized space for information dissemination, to ensure that any relevant parties can stay informed of any news or announcement. Second goal is to create an easy-to-access educational digital discussion platform in order to foster communication and discussions among students university staff. 

Technically, the website was built using MERN stack (MongoDB, Express.js, React.js, and Node.js), which is a website development framework that is famous for its quick development time, ease of use, and its compatibility with JavaScript language, with the framework being complete on its own, covering front end, back end, API, and database, all within one framework.

The research concludes that transitioning to a digital framework significantly improves user engagement and reduces the risk of miscommunication. For future development, the author suggests further optimizing the user experience and exploring additional features.

\subsection{Course Selling Website Using MERN Stack \cite{paper_12}}

This paper presents an educational and commercial website that lets university teacher to create, manage, and sell their courses online, intended to be a platform that improves the education quality by narrowing the gap between student's demand for educational courses and teacher's courses supply under the platform equipped with firm security and secure user's privacy protection.

With the website development heavily implements the MERN stack (MongoDB, Express.js, React.js, Node.js) which is a website development framework that is known for its quick development time, ease of use, and its high compatibility with JavaScript language. Additionally, the website also implements the REST API to obtain website's fast response time, and the capability to handle a large number of users and website's course data without being slowed down.

The study found that this architecture makes the platform very flexible, allowing for features like real-time chat and secure payments to be added easily. By keeping the "business side" of the website separate from the "user side," the system is easier to update as more students join.

For future work, the author suggests adding mobile options for the website, AI assistance to help answer student's questions, and digital certificate to serve as a proof when students finish their courses.


\section{Functions}

\funcgroup{Schedule matcher}

\subsection{OCR with Tesseract, Amazon Textract, and Google Document AI: a benchmarking experiment \cite{paper_16}}
This study addresses the variability in OCR software accuracy by conducting a benchmarking experiment comparing the performance of Tesseract, Amazon Textract, and Google Document AI on images of English and Arabic text. 

English-language book scans (n=322) and Arabic-language article scans (n=100) were replicated 43 times with different types of artificial noise for a corpus of 18,568 documents, generating 51,304 process requests. 

Performance was evaluated by measuring character and word error rates across document types and noise conditions, comparing all three engines under identical test materials and standardised noise injection.

The study demonstrated that Document AI and Textract performed significantly better than Tesseract, particularly on noisy documents, while Tesseract remains a cost-effective option but struggles with complex layouts. 

Future work includes testing additional OCR processors such as Adobe PDF Services and ABBYY Cloud OCR using the same benchmark procedure, and extending the dataset to more language families beyond English and Arabic.

\subsection{Printed document layout analysis and optical character recognition system based on deep learning \cite{paper_17}}
This study addresses the challenge of automated document understanding by proposing a layout analysis and text recognition system for printed documents based on deep learning. 

Scanned documents or image files are processed using a layout analysis algorithm based on YOLOv4 and YOLOv8 deep learning to identify the positions of titles, text paragraphs, tables, and images within the document, after which each category undergoes specific character segmentation processing before content is recognised using a CNN-based text recognition algorithm. 

Performance was evaluated on printed document datasets measuring layout detection accuracy and character recognition rates across different document structures and categories.

The proposed system demonstrates strong performance on structured printed documents by combining layout detection and text recognition into an end-to-end pipeline, with the layout analysis stage accurately separating structural elements before passing them to the recognition stage. 

Future work includes extending the system to handle handwritten documents and more complex multi-column layouts, as well as improving robustness to degraded or low-resolution scans.

\subsection{Color sensors and their applications based on real-time color image segmentation for cyber physical systems \cite{paper_18}}
This study addresses the challenge of real-time color detection and segmentation in cyber physical systems by proposing a novel real-time color image segmentation method based on color similarity in RGB color space. 

The proposed method uses OpenCV's CVCAM technology to capture and process video streams, determining the dominant color first based on RGB color and luminance information, then calculating color similarity using a proposed color component method to create a color-class map. 

Performance was evaluated by applying the segmentation algorithm to real-time video streams and measuring the accuracy of color region detection under varying lighting and background conditions.

The experiment demonstrated that the proposed RGB-based segmentation method can reliably identify and isolate colored regions in real time, with results showing precise segmentation that handles noise and lighting variation effectively. 

Future work includes extending the method to more complex multi-object scenes and integrating adaptive color correction mechanisms to further improve robustness under extreme lighting conditions.

\subsection{A hyper-heuristic algorithm based on genetic and greedy strategy for university course scheduling problem \cite{paper_3}}

This study addresses the University Course Scheduling Problem (UCSP) by proposing a novel model that balances resources like classrooms and teachers while satisfying complex constraints.

A hyper-heuristic algorithm (GA\_RG\_HH) was developed, combining Genetic Algorithms (GA) with greedy strategies to optimize scheduling efficiency.
Performance was evaluated using real-world datasets, measuring the re- duction of soft constraint violations and the efficiency of resource allocation.

The experiment indicated that the proposed GA\_RG\_HH performs bet- ter than traditional algorithms like Tabu Search and Particle Swarm Optimization. This is attributed to the algorithm's ability to balance exploration and exploitation through a pre-check strategy and dynamic operator sequences.
Future work includes further optimization of the algorithm for even larger-scale institutional data and the inclusion of more diverse constraint parameters.

\subsection{Algorithm Design \cite{kleinberg2005algorithm}}
This chapter addresses the Interval Scheduling Problem, which asks how to select the maximum number of non-conflicting tasks from a set of requests, each defined by a start time and a finish time, such that no two selected tasks overlap in time.

A greedy algorithm was proposed, selecting tasks by earliest finish time at each step — that is, always choosing the next available task that ends the soonest — and proving that this simple strategy consistently produces an optimal solution without requiring exhaustive search.
Performance was analyzed theoretically using an exchange argument, demonstrating that the greedy solution is always at least as good as any other feasible selection, and the algorithm was shown to run in O(n log n) time after sorting tasks by finish time.

The result established that the earliest-finish-time greedy strategy is provably optimal for the basic interval scheduling problem, outperforming naive approaches such as selecting by shortest duration or earliest start time, which can both produce suboptimal results.

The chapter further extends the framework to related problems including Interval Partitioning, which finds the minimum number of resources needed to schedule all tasks without conflict, and Scheduling to Minimize Lateness, laying groundwork for more complex scheduling variants in future study.

\funcgroup{Course rater}

\subsection{A Review of Strategies for Designing, Administering, and Using Student Ratings of Instruction \cite{paper_14}}

This paper reviews the best ways to design and use student feedback (course evaluations) to help teachers improve their classes, specifically in medical and pharmacy programs.

The core of the study suggests that a good evaluation should be carefully planned to be fair and useful. The author recommended using a standardized online forms with about 10 to 30 questions that are easy to understand for students. The reason why the forms are suggested to be in online format is because it allows the results to be collected and processed more efficiently compared to the traditional paper format. The research also recommended that there should be a formal scheduled time for students to complete the survey, for example, during class period, to prevent "survey fatigue".

Additionally, the research found that the evaluations are most effective when they focus on the "teaching aspect" of the courses, for example, the goal of the course, and how well the the teachers prepare and present their materials. The study also pointed out that although online evaluations receive less response than paper evaluations, the quality of the results is usually the same.

For future work, the author suggests using random samples with the survey, since the traditional method of collecting course evaluations is to distribute the survey to every students, but by using random sampling (selecting a group of students in random), it can reduce the time taken to collect the evaluations while maintaining the quality of the results. The teacher's access to the evaluation data is also advised to be limited for better efficiency in data distribution. 

\funcgroup{Super planner}

\subsection{CourseViewer a prototype for visualizing undergraduate courses and student transcripts \cite {paper_15}}

This paper introduces CourseViewer, a prototype that transforms traditional, text-heavy transcripts into interactive visual maps to help students plan their academic paths. This approach is highly relevant for an academic planner because it replaces manual checking with a dynamic system that shows how different courses are connected.

The most useful feature for a modern planner is the graph-based visualization. By representing courses as "nodes" and prerequisites as directed arrows, the system makes it easy to see which classes "unlock" future subjects. This technical design helps prevent enrollment errors by clearly showing if a student has met the necessary requirements before they try to add a class.

Technically, the tool uses the Java language and the Prefuse visualization toolkit to automatically handle the map layout so the user doesn't have to organize it manually. It also uses color-coding—such as green for completed and red for required—which provides an immediate visual status of a student's progress toward graduation.

The study found that these visual maps allow students to identify their remaining credits much faster and more accurately than reading a paper transcript. For future work, the authors suggested adding a "what-if" feature to simulate different majors and integrating the tool directly into online registration systems.


\subsection{The curriculum prerequisite network: a tool for visualizing and analyzing academic curricula \cite{paper_13}}

This paper treats a university curriculum as a "complex system" that can be mapped using network science to reveal the hidden connections between courses. This is directly applicable to the project’s goal of automating 4-year academic planning for major and general education courses, which currently must be done manually.

The study uses a Directed Acyclic Graph (DAG) where courses are "nodes" and prerequisites are "links". A key technical detail is the identification of "hubs" (courses that unlock many others) and "bridges" (links between different course groups). Implementing this logic allows the interactive planner to visually flag "bottleneck" courses, helping students avoid choices that could delay their graduation by a year.

Technically, the research utilized NetworkX and Gephi to analyze network density and "betweenness centrality". 

\section{Contribution}

The literature reviewed in this chapter provides both technical and conceptual foundations that directly inform the design and development of EGCEye.

With the three papers from technological stack section of this chapter, both \cite{paper_11} and \cite{paper_12} built their platforms using MongoDB, Express.js, React, and Node.js, and both worked well for educational use cases at a scale close to EGCEye. So it's fair to say the MERN stack is a solid, well-tested choice for this kind of project. And the REST API approach from \cite{paper_10} shows that one API layer can handle multiple platforms at the same time without things getting out of sync, which is exactly how EGCEye is set up too, so it is fair to say that the majority of the contributions of the three papers is they provide a proven-to-work examples of website development using MERN stack which serves as a suggestion of our website development planning.

For the Schedule Matcher, Hegghammer \cite{paper_16} provides peer-reviewed justification for the selection of Tesseract as the OCR engine by benchmarking it against commercial alternatives, confirming it as a cost-effective choice for clean, structured document inputs such as Sky+ schedule screenshots. Li et al. \cite{paper_17} further supports the two-stage pipeline design of first detecting the schedule grid layout before applying OCR, validating that separating layout detection from text recognition improves overall accuracy. Xiong et al. \cite{paper_18} establishes the theoretical basis for the color blob detection step, demonstrating that OpenCV-based RGB color segmentation reliably isolates colored regions against high-contrast backgrounds. Finally, Kleinberg and Tardos \cite{kleinberg2005algorithm} provide the formal interval scheduling theory that underpins both the busy map OR logic and the subgroup fallback algorithm.

For the Rate-My-Course function, \cite{paper_14} provides evidence-based guidance on feedback system design, supporting EGCEye's decision to keep star rating as the only required field, list courses separately by professor where multiple sections exist, and present aggregated results in a structured and accessible format.

For the Super Planner, CourseViewer \cite{paper_15} serves as the direct precedent for the node-and-arrow graph representation and the color-coding of course completion states. The curriculum prerequisite network study \cite{paper_13} further supports the use of graph-based visualization over text-based planning tools, and informs the planner's bottleneck warning system that flags courses a student should prioritize based on their year level.
